\section{Kultureller und technischer Kontext}

Auf der Suche nach einer unmittelbaren räumlichen Realität experimentierte der italienische Künstler Duccio di Buoninsegna (ca. 1255–1319) mit der perspektivischen Darstellung. In seinem Werk \textit{Das Letzte Abendmahl} (ca. 1310) zeichnete er erste Linienpaare mit einem eigenen Fluchtpunkt, während in dem späteren Werk \textit{Ankündigung des Todes der Jungfrau} (1310) schließlich alle Deckenlinien in einem Hauptfluchtpunkt zusammenlaufen. \footnote{Vgl.: \cite[S. 4-9]{Andersen2007}}

Als Begründer der eigentlichen Perspektive im zweiten Jahrzehnt des \textit{quattrocento} gilt jedoch Filippo Brunelleschi (1377–1446). Laut frühen Biografien demonstrierte er seine Erfindung auf zwei Holztafeln, welche die Kirche \textit{San Giovanni} und den \textit{Palazzo dei Signori} in Florenz zeigten. Für die erste Holztafel nutzte Brunelleschi einen Spiegel, über den der Beobachter auf das Spiegelbild des Baptisteriums blickte, was den Anspruch untermauerte, eine exakte Momentaufnahme der Wirklichkeit zu reproduzieren. Trotz dieser historischen Bedeutung bleibt bis heute unklar, woher er die Inspiration für seinen Aufbau fand oder welche mathematische Konstruktionsmethode er für die Holztafeln anwandte, da die Originale verloren gingen. Ausgehend von einer Gruppe Florentiner Künstler verbreitete sich Brunelleschis Innovation jedoch über ganz Italien, sodass die Zentralperspektive bis zum Ende des 16. Jahrhunderts zu einem allgemeingültigen Paradigma für die folgenden 300 Jahre avancierte. \footnote{Vgl.: \cite[S. 11-14]{Andersen2007}}

Die perspektivische Anamorphose als gezielte zeichnerische Verzerrung entsprach in der Folge exakt der Vorliebe des Manierismus für das Mysteriöse. \footnote{Vgl.: \cite[S. 71-75]{Andersen2007}} In dieser Epoche erlangten Künstler wie Giuseppe Arcimboldo (1527-1593) durch Werke wie \textit{Vertumnus} (1590) Berühmtheit, indem sie Porträts rein aus Objekten wie Früchten oder Blumen komponierten. \footnote{Vgl.: \cite{Kemp1990}} Während solche manieristischen Werke oft auf einer spielerischen Verzerrung durch einfache Rastermethoden beruhten, unterschied sich dies deutlich von der mathematischen Präzision der direkten Konstruktion nach Piero della Francesca. \footnote{Vgl.: \cite[S. 74-75]{Andersen2007}}

Ihren Höhepunkt fand diese Entwicklung schließlich in den Arbeiten des Minim-Priesters Jean-François Niceron sowie des Jesuiten Athanasius Kircher und dessen Schüler Gaspar Schott. Sie nutzten diese optischen Wunderwerke im Kontext der „natürlichen Magie“ und der Illusionskunst primär für didaktische und religiöse Zwecke, um Wissen auf faszinierende Weise zu vermitteln. \footnote{Vgl.: \cite[S. 609-611]{Andersen2007}}